For each \(m\), let \(\mathcal T_m\) be a finite index set and let
\((Y_{m,T})_{T\in\mathcal T_m}\) be nonnegative random variables.  Put
\[
  Y_m^*=\min_{T\in\mathcal T_m}Y_{m,T}.
\]

\begin{definition}[Exponential-scale small-ball data]
Fix \(H>0\) and \(a>0\).  We say that the family has upper small-ball exponent
\(a\) and entropy \(H\) if
\[
  \log |\mathcal T_m|\le (H+o(1))m
\]
and, for every fixed \(s>0\),
\[
  \sup_{T\in\mathcal T_m}
  \PP\{Y_{m,T}\le e^{-sm}\}
  \le \exp\{-(as+o(1))m\}.
\]
It has a second-moment lower small-ball principle below \(H/a\) if, for every
fixed \(s<H/a\), the counting variable
\[
  Z_m(s)=\sum_{T\in\mathcal T_m}\mathbf{1}_{\{Y_{m,T}\le e^{-sm}\}}
\]
satisfies
\[
  \EE Z_m(s)\to\infty,
  \qquad
  \frac{\Var Z_m(s)}{(\EE Z_m(s))^2}\to0.
\]
\end{definition}

\begin{theorem}[Extremal small-ball principle]\label{thm:principle}
Assume the upper small-ball exponent and entropy condition above.  Then, for
every \(\varepsilon>0\),
\[
  \PP\left\{Y_m^*\le \exp\left[-\left(\frac{H}{a}+\varepsilon\right)m\right]\right\}
  \longrightarrow 0.
\]
If, in addition, the second-moment lower small-ball principle holds below
\(H/a\), then
\[
  \frac{1}{m}\log Y_m^*\xrightarrow{\PP}-\frac{H}{a}.
\]
\end{theorem}

\begin{proof}
Let \(s>H/a\).  By the union bound,
\[
  \PP\{Y_m^*\le e^{-sm}\}
  \le |\mathcal T_m|
      \sup_{T\in\mathcal T_m}\PP\{Y_{m,T}\le e^{-sm}\}
  \le \exp\{(H-as+o(1))m\},
\]
which tends to zero.  This proves the lower bound on \(Y_m^*\).

Let \(s<H/a\).  If the second-moment hypothesis holds, then Chebyshev gives
\[
  \PP\{Z_m(s)=0\}
  \le \frac{\Var Z_m(s)}{(\EE Z_m(s))^2}\to0.
\]
Thus, with probability tending to one, some \(T\) has \(Y_{m,T}\le e^{-sm}\).
This says
\[
  \PP\left\{Y_m^*\le \exp\left[-\left(\frac{H}{a}-\varepsilon\right)m\right]\right\}
  \to1
\]
after taking \(s=H/a-\varepsilon\).  The two estimates are exactly convergence
in probability of \(m^{-1}\log Y_m^*\) to \(-H/a\).
\end{proof}

\begin{corollary}[Polynomial small-ball version]\label{cor:poly-smallball}
Suppose
\[
  |\mathcal T_m|\le \exp((H+o(1))m)
\]
and for some constants \(C,B,a>0\),
\[
  \sup_T\PP\{Y_{m,T}\le t\}\le C m^B t^a,\qquad 0<t\le1.
\]
Then for every \(\varepsilon>0\),
\[
  \PP\left\{Y_m^*\le \exp\left[-\left(\frac{H}{a}+\varepsilon\right)m\right]\right\}
  \to0.
\]
If the matching second-moment condition holds below \(H/a\), then
\((1/m)\log Y_m^*\to-H/a\) in probability.
\end{corollary}

\begin{corollary}[Finite-budget lower side]\label{cor:finite-budget}
Let \(\mathcal T_m\) be finite and suppose that, for some \(C_m\ge1\) and
\(a_m>0\),
\[
  \sup_{T\in\mathcal T_m}\PP\{Y_{m,T}\le t\}\le C_m t^{a_m},
  \qquad 0<t\le1.
\]
Then, for every \(u_m>0\),
\[
  \PP\left\{
    Y_m^*\le
    \exp\left[-\frac{\log|\mathcal T_m|+\log C_m+u_m}{a_m}\right]
  \right\}
  \le e^{-u_m}.
\]
\end{corollary}

\begin{proof}
Put \(t_m=\exp[-(\log|\mathcal T_m|+\log C_m+u_m)/a_m]\).  Since
\(|\mathcal T_m|,C_m\ge1\), one has \(0<t_m\le1\).  The union bound gives
\[
  \PP\{Y_m^*\le t_m\}
  \le |\mathcal T_m| C_m t_m^{a_m}
  = e^{-u_m}.
\]
\end{proof}

\begin{corollary}[Nonuniform finite-budget lower side]\label{cor:nonuniform-finite-budget}
Let \(\mathcal T_m\) be nonempty and finite.  Suppose that, for deterministic
numbers \(a_m>0\) and \(C_{m,T}\ge0\),
\[
  \PP\{Y_{m,T}\le t\}\le C_{m,T}t^{a_m},
  \qquad 0<t\le1,\quad T\in\mathcal T_m.
\]
If \(B_m\ge1\vee\sum_{T\in\mathcal T_m}C_{m,T}\), then, for every \(u_m>0\),
\[
  \PP\left\{
    Y_m^*\le \exp\left[-\frac{\log B_m+u_m}{a_m}\right]
  \right\}
  \le e^{-u_m}.
\]
\end{corollary}

\begin{proof}
The displayed threshold is at most one.  Hence the union bound gives
\[
  \PP\{Y_m^*\le t\}
  \le \sum_{T\in\mathcal T_m}C_{m,T}t^{a_m}
  \le B_m t^{a_m},\qquad 0<t\le1.
\]
Substitute \(t=\exp[-(\log B_m+u_m)/a_m]\).
\end{proof}

\begin{corollary}[Almost-sure nonuniform finite-budget lower side]\label{cor:as-nonuniform-finite-budget}
Assume that the random variables for all \(m\) are realized on a common
probability space.  Suppose that the index sets \(\mathcal T_m\) are nonempty,
finite, and deterministic, and that \(a_m>0\), \(u_m>0\),
\(C_{m,T}\ge0\), and \(B_m\ge1\vee\sum_{T\in\mathcal T_m}C_{m,T}\) are
deterministic.  If
\[
  \PP\{Y_{m,T}\le t\}\le C_{m,T}t^{a_m},
  \qquad 0<t\le1,\quad T\in\mathcal T_m,
\]
and
\[
  \sum_{m=1}^\infty e^{-u_m}<\infty,
\]
then, almost surely, for all sufficiently large \(m\),
\[
  Y_m^*>
  \exp\left[-\frac{\log B_m+u_m}{a_m}\right].
\]
\end{corollary}

\begin{proof}
Apply Corollary~\ref{cor:nonuniform-finite-budget} at the displayed threshold
and then use the first Borel--Cantelli lemma.  No independence between
different values of \(m\), or between the events inside one \(\mathcal T_m\),
is used.
\end{proof}

\begin{corollary}[Almost-sure finite-budget lower side]\label{cor:as-finite-budget}
Assume that the random variables for all \(m\) are realized on a common
probability space.  Suppose that the index sets \(\mathcal T_m\) are nonempty,
finite, and deterministic, and that the numbers \(C_m\ge1\), \(a_m>0\), and
\(u_m>0\) are deterministic.  If
\[
  \sup_{T\in\mathcal T_m}\PP\{Y_{m,T}\le t\}\le C_m t^{a_m},
  \qquad 0<t\le1,
\]
and
\[
  \sum_{m=1}^\infty e^{-u_m}<\infty,
\]
then, almost surely, for all sufficiently large \(m\),
\[
  Y_m^*>
  \exp\left[-\frac{\log|\mathcal T_m|+\log C_m+u_m}{a_m}\right].
\]
\end{corollary}

\begin{proof}
Put
\[
  t_m=\exp\left[-\frac{\log|\mathcal T_m|+\log C_m+u_m}{a_m}\right].
\]
Then \(0<t_m\le1\), and the proof of Corollary~\ref{cor:finite-budget} gives
\[
  \PP\{Y_m^*\le t_m\}\le e^{-u_m}.
\]
The first Borel--Cantelli lemma gives
\(\PP\{Y_m^*\le t_m\ {\rm infinitely\ often}\}=0\).  Hence
\(Y_m^*>t_m\) eventually almost surely.  No independence between different
values of \(m\), or between the events inside one \(\mathcal T_m\), is used.
\end{proof}

\begin{corollary}[Almost-sure lower envelope]\label{cor:as-lower-envelope}
Assume that the random variables for all \(m\) are realized on a common
probability space.  Suppose there is a constant \(c\ge0\) such that, for every
\(\varepsilon>0\),
\[
  \sum_{m=1}^\infty
  |\mathcal T_m|\,
  \sup_{T\in\mathcal T_m}
  \PP\{Y_{m,T}\le e^{-(c+\varepsilon)m}\}<\infty .
\]
Then
\[
  \liminf_{m\to\infty}\frac1m\log Y_m^*\ge -c
  \qquad\hbox{almost surely}.
\]
In particular, the upper small-ball exponent and entropy hypothesis in
Theorem~\ref{thm:principle} imply the almost-sure lower envelope with
\(c=H/a\).
\end{corollary}

\begin{proof}
For fixed \(\varepsilon>0\), the event
\[
  E_m(\varepsilon)=\{Y_m^*\le e^{-(c+\varepsilon)m}\}
\]
has probability bounded by the summand in the hypothesis.  The first
Borel--Cantelli lemma gives \(\PP\{E_m(\varepsilon)\ \hbox{infinitely often}\}=0\).
Intersecting the resulting probability-one events over
\(\varepsilon=1,1/2,1/3,\ldots\) gives the displayed liminf.  No independence
between different \(m\) is used.

For the last assertion, take \(s=H/a+\varepsilon\).  The union bound in
Theorem~\ref{thm:principle} gives
\[
  \PP\{Y_m^*\le e^{-sm}\}\le \exp\{(-a\varepsilon+o(1))m\},
\]
which is summable.
\end{proof}

The Gaussian source theorem is the case
\[
  \mathcal T_m=\{T\subset[N_m]: |T|=m\},\qquad
  H=h(\gamma),\qquad
  a=d_\F,
\]
with the second-moment condition proved in \cite{ShmaloExtremal2026}.
